\section{Literature Review}
\label{sec:literature}

\begin{table*}[!t]
\centering
\caption{Summary of representative prior studies --- dataset, models, performance, and experimental gaps.}
\label{tab:prior_work}
\scriptsize
\setlength{\tabcolsep}{4pt}
\begin{tabular}{@{}p{3.2cm}p{3.2cm}p{2.8cm}cp{4.2cm}@{}}
\toprule
\textbf{Study} & \textbf{Dataset} & \textbf{Models} & \textbf{Best F1} & \textbf{Key Limitation} \\
\midrule
\citet{saravia2018carer} & Twitter (6 cls, $\sim$20K) & CNN + word emb. & 0.79 macro & Small dataset; single model family \\
\citet{hasan2019emotion} & Twitter multi-class & CNN, LSTM & $\sim$0.82 acc. & No class-balancing; accuracy only \\
\citet{demszky2020goemotions} & GoEmotions --- Reddit (27 cls, 58K) & BERT fine-tuned & 0.46 macro & Different domain; 27 classes \\
\citet{sun2023bidirectional} & Twitter sentiment & BiLSTM & $\sim$0.85 acc. & Accuracy-only; inconsistent preproc. \\
\citet{acheampong2021transformer} & Multiple benchmarks & BERT (survey) & N/A & Survey only; no unified comparison \\
\midrule
\textbf{This Study} & \textbf{Kaggle Twitter (6 cls, 416K)} & \textbf{SVM, XGBoost, CNN, BiLSTM} & \textbf{0.94} & \textbf{Unified pipeline; macro-F1 eval.} \\
\bottomrule
\end{tabular}
\end{table*}

Emotion detection research has progressed from lexicon-based methods~\citep{mohammad2013crowdsourcing, strapparava2008learning} to machine learning and deep learning architectures. Traditional ML models such as SVM and Na\"ive Bayes perform well on structured data with TF-IDF features but cannot capture sequential context~\citep{dzikovska2012using, wang2018emotion}. Deep learning methods --- including CNNs for local n-gram feature extraction~\citep{hasan2019emotion} and BiLSTMs for long-range sequential dependencies~\citep{sun2023bidirectional} --- learn contextual patterns directly from text. XGBoost provides a strong non-linear baseline through gradient boosting with L1/L2 regularisation~\citep{chen2016xgboost}. More recently, transformer-based models such as BERT achieve state-of-the-art results through contextual embeddings, though at substantially higher computational cost~\citep{acheampong2021transformer, devlin2019bert}. Across these approaches, class imbalance and inconsistent evaluation remain persistent challenges~\citep{nandwani2021review}.

Table~\ref{tab:prior_work} summarises representative studies, highlighting datasets, models, performance, and experimental gaps. To the best of our knowledge, no prior work evaluates SVM, XGBoost, CNN, and BiLSTM on the same large-scale Twitter corpus under identical preprocessing, stratified splits, and macro-F1 evaluation. This study fills that gap through a controlled comparison of classical and sequence-based architectures under a unified experimental framework.
